\chapter{Diseño experimental}
\label{chap:diseno}

Este trabajo parte de una limitación concreta: el número de pacientes etiquetados es reducido y la entrada disponible es una imagen del ECG. El diseño experimental compara dos formas de aprovechar esos ejemplos. Ambas emplean información visual, mantienen la separación entre pacientes y se evalúan con las mismas métricas.

En el enfoque supervisado, una ResNet18 aprende a partir de los ejemplos disponibles y actualiza sus pesos durante el entrenamiento. Gemma~4, en cambio, recibe los ejemplos como demostraciones dentro del prompt, sin modificar sus parámetros. Esta distinción determina el significado de \(K\): en la CNN denota pacientes de entrenamiento y, en el VLM, demostraciones incluidas en el contexto.

El sistema se plantea como una herramienta experimental de priorización, no como un método diagnóstico. Una predicción positiva identifica un trazado que convendría revisar, pero no equivale a confirmar ni a descartar el síndrome de Brugada. La valoración clínica requeriría el ECG completo y, según el caso, otras pruebas como el reposicionamiento de V1/V2 o el test farmacológico. Desde esta perspectiva, la sensibilidad indica qué proporción de positivos genera una alerta, mientras que la especificidad refleja cuántos pacientes negativos no serían derivados innecesariamente.

El protocolo distingue los datos, las particiones, las métricas, el papel previsto del sistema y los resultados. Además, incluye matrices \(2\times2\), sensibilidad y especificidad para facilitar la interpretación \cite{collins2024tripodai,bossuyt2015stard,sounderajah2025stardai}. Al trabajar con imágenes clínicas, también se especifican la entrada del modelo, la agregación de las predicciones y la información que queda fuera de la imagen. La Figura~\ref{fig:diseno_experimental_resumen} resume las ramas comparadas.

\begin{figure}[H]
  \centering
  \safeincludegraphics[width=0.96\textwidth]{results/fig_pipeline_experimental.drawio.png}
  \caption{Resumen del diseño experimental común: entrada de ECG como imagen, uso supervisado de ejemplos en la CNN, uso contextual de ejemplos en el VLM/ICL y evaluación con métricas compartidas.}
  \label{fig:diseno_experimental_resumen}
\end{figure}

\section{Planteamiento común}

La rama supervisada entrena una ResNet18 con pocos ejemplos. La rama VLM/ICL utiliza Gemma~4 y contempla tres condiciones: zero-shot, ICL estándar e ICL balanceado. El uso de los mismos valores de \(K\), la separación por paciente y un procedimiento común de agregación permite compararlas bajo un protocolo homogéneo, aunque cada modelo aprovecha los ejemplos etiquetados de manera distinta.

La comparación abarca dos dominios. El simulador QRS/ST permite estudiar si las imágenes contienen información morfológica discriminativa en un entorno controlado. La cohorte del HUCA introduce ECG clínicos reales, con ruido de adquisición, variabilidad entre pacientes y una etiqueta cuyo significado no coincide exactamente con las reglas empleadas en la simulación.

\section{Datos evaluados}

El paciente constituye siempre la unidad de separación. En el conjunto sintético, cada paciente simulado aporta una imagen de V1. En la cohorte del HUCA se utiliza, de forma análoga, una única ventana renderizada de V1 por paciente. Cada imagen genera una predicción, que se interpreta directamente a nivel de paciente.

\begin{table}[H]
  \centering
  \small
  \begin{tabularx}{\textwidth}{p{3cm}rrrrX}
    \toprule
    Conjunto & Pacientes & Imágenes & Positivos & Negativos & Tarea \\
    \midrule
    Sintético QRS/ST & 100 & 100 & 20 & 80 & Patrón derivado de Brugada con RBBB, elevación ST e inversión T presentes \\
    Datos HUCA & 317 & 317 & 116 & 201 & Referencia binaria de la campaña final tras el filtrado local \\
    \bottomrule
  \end{tabularx}
  \caption{Conjuntos utilizados en los experimentos finales. Cada paciente aporta una imagen; las matrices de confusión acumulan las predicciones obtenidas con las tres semillas.}
  \label{tab:datasets_finales}
\end{table}

Las etiquetas no representan exactamente el mismo fenómeno en ambos dominios. En el simulador, la clase positiva responde a una regla morfológica definida mediante criterios explícitos. En la cohorte del HUCA, procede de una variable clínica incluida en los metadatos. Por tanto, el paso del entorno sintético al clínico introduce tanto un cambio visual como un cambio semántico.

\section{Valores de K y validación}

Los experimentos con pocos ejemplos emplean \(K=\{2,4,8,16,32\}\). En el VLM se añade \(K=0\), correspondiente a la condición zero-shot. Cada valor positivo se evalúa con las semillas 42, 123 y 2026. Los ejemplos se seleccionan de nuevo en cada pliegue, sin incluir nunca al paciente reservado para prueba.

Las variantes de ICL se diferencian únicamente en esa selección. Para \(K\geq2\), el ICL estándar incorpora al menos una demostración positiva y otra negativa; las posiciones restantes se completan de forma seudoaleatoria entre los pacientes disponibles. No se trata, por tanto, de una recuperación aleatoria pura. El ICL balanceado intenta asignar \(K/2\) demostraciones a cada clase y solo completa con la clase disponible si una de ellas no reúne suficientes casos. En ambos métodos, la semilla fija la selección y el paciente de prueba queda excluido.

En la evaluación sobre la cohorte del HUCA, el origen de \(K\) varía entre las ramas. La CNN se ajusta con \(K\) pacientes del simulador, ya que la cohorte clínica no contiene las etiquetas morfológicas requeridas por la cabeza multietiqueta. En el VLM, \(K\) representa demostraciones de contexto. Para la comparación clínica principal se toman de la propia cohorte; en el análisis de transferencia, proceden del simulador. La Figura~\ref{fig:k_examples_split} muestra la separación por paciente.

\begin{figure}[H]
  \centering
  \safeincludegraphics[width=0.94\textwidth]{results/fig_k_examples_split.pdf}
  \caption{Separación por paciente y selección de ejemplos para cada valor de \(K\). El paciente reservado no interviene en la selección; los ejemplos actualizan los pesos de la CNN o forman el contexto del VLM.}
  \label{fig:k_examples_split}
\end{figure}

Esta separación evita que registros de una misma persona aparezcan a la vez en entrenamiento y prueba, una fuente habitual de sesgo en el análisis de ECG. Sin este control, el modelo podría reconocer rasgos del paciente, del dispositivo o del renderizado en vez de patrones relacionados con la tarea.

\section{Métricas}

La métrica principal es la \textit{Balanced Accuracy} (BA), que asigna el mismo peso a la sensibilidad y a la especificidad aunque las clases estén desbalanceadas. Se calcula como

\[
  \mathrm{BA}=\frac{1}{2}\left(\frac{TP}{TP+FN}+\frac{TN}{TN+FP}\right).
\]

También se informan F1, sensibilidad, especificidad, área bajo la curva ROC y precisión media. Las tablas principales presentan la media de las tres semillas y, cuando procede, su desviación estándar. Las matrices de confusión acumulan las predicciones de las tres ejecuciones.

Una sensibilidad alta acompañada de una especificidad baja puede resultar útil en un análisis exploratorio, porque permite recuperar más positivos. Sin embargo, no sería suficiente para un circuito asistencial si generase demasiadas alertas. Esta tensión orienta la lectura de las comparaciones posteriores.

El capítulo~\ref{chap:modelos} desarrolla la CNN y el enfoque VLM/ICL dentro de este marco comparativo. Después, el capítulo~\ref{chap:resultados} presenta los resultados cuantitativos obtenidos en el dominio sintético y en la cohorte del HUCA.
